% Gerado por scripts/migrate_markdown_to_latex.py.
% Fonte migrada: fases/01_validacao_conceitual/evidencias/benchmarks/diagnostico_cpu_l64_d128_2026-08-25/analise/relatorio_preliminar.md


\chapter{Relatório preliminar - diagnostico\_\allowbreak{}cpu\_\allowbreak{}l64\_\allowbreak{}d128\_\allowbreak{}v1}



\begin{icquote}
Esta coleta é um diagnóstico em CPU de um ponto da grade formal. Os tempos não substituem a coleta na GPU-alvo e não sustentam conclusão de hardware.
\end{icquote}



\section{Pergunta, hipótese e desenho}


\begin{itemize}
\item Pergunta principal: quantização INT8 simulada preserva melhor a saída que pruning por magnitude de 50\% nas mesmas operações?
\item Hipótese: a quantização terá MSE/MAE menores e R²/cosseno maiores que o pruning.
\item Amostra: tensores sintéticos de distribuição normal, gerados deterministicamente; não há dataset ou pré-processamento externo.
\item Seed: \texttt{42}; lote(s): \texttt{[1]}; sequências: \texttt{[64]}; dimensões: \texttt{[128]}.
\item Medição: \texttt{20} warm-ups, \texttt{50} medições e \texttt{2} execuções independentes.
\item Dispositivo observado: \texttt{cpu}
\item Inclusão: todos os casos válidos da grade e os três cenários registrados; exclusão: qualquer artefato com checksum, shape, dtype ou valor numérico inválido.
\end{itemize}


\section{Reprodutibilidade}


\begin{itemize}
\item Hashes de entrada iguais entre execuções: \texttt{True}.
\item Hashes de saída iguais entre execuções: \texttt{True}.
\item As latências não precisam ser idênticas: elas medem ruído e estado do sistema; entradas e saídas determinísticas precisam coincidir.
\end{itemize}


\section{Resultados por operação e cenário}



\begin{landscape}
\scriptsize
\begin{longtable}{@{}>{\RaggedRight\arraybackslash}p{0.117\linewidth}>{\RaggedRight\arraybackslash}p{0.117\linewidth}>{\RaggedRight\arraybackslash}p{0.117\linewidth}>{\RaggedRight\arraybackslash}p{0.117\linewidth}>{\RaggedRight\arraybackslash}p{0.117\linewidth}>{\RaggedRight\arraybackslash}p{0.117\linewidth}>{\RaggedRight\arraybackslash}p{0.117\linewidth}@{}}
\toprule
Operação & Cenário & Latência média (ms) & DP entre execuções & MSE & R² & Cosseno \\
\midrule
\endfirsthead
\toprule
Operação & Cenário & Latência média (ms) & DP entre execuções & MSE & R² & Cosseno \\
\midrule
\endhead
dense\_\allowbreak{}projection & baseline & 0.088291 & 0.0207253 & 0 & 1 & 1 \\
dense\_\allowbreak{}projection & pruning\_\allowbreak{}magnitude & 0.076017 & 0.0643283 & 9.30825 & 0.927085 & 0.962897 \\
dense\_\allowbreak{}projection & quantization\_\allowbreak{}int8 & 0.061625 & 0.00428648 & 0.00995697 & 0.999922 & 0.999961 \\
self\_\allowbreak{}attention & baseline & 0.93542 & 0.383818 & 0 & 1 & 1 \\
self\_\allowbreak{}attention & pruning\_\allowbreak{}magnitude & 0.609696 & 0.0349735 & 12144.4 & 0.242143 & 0.596755 \\
self\_\allowbreak{}attention & quantization\_\allowbreak{}int8 & 0.867666 & 0.183593 & 732.65 & 0.95428 & 0.977141 \\
\bottomrule
\end{longtable}
\end{landscape}



\section{Interpretação preliminar}


\begin{itemize}
\item \texttt{dense\_\allowbreak{}projection}: MSE da quantização \texttt{0.\allowbreak{}00995697} e do pruning \texttt{9.\allowbreak{}30825}; evidência compatível com H1 nesta amostra.
\item \texttt{self\_\allowbreak{}attention}: MSE da quantização \texttt{732.\allowbreak{}65} e do pruning \texttt{12144.\allowbreak{}4}; evidência compatível com H1 nesta amostra.
\item Pruning usa matriz densa e não usa kernel esparso; quantização é dequantizada para \texttt{float32} e não usa kernel INT8.
\item Razões de latência e memória estão nos CSVs de análise, mas só podem sustentar alegação de hardware se o experimento for principal e o caminho executado tiver validade \texttt{hardware}.
\item Resultados de smoke ou diagnóstico CPU não são estimativas finais de desempenho no hardware-alvo.
\end{itemize}


\section{Cadeia de evidência}



\texttt{pergunta → H1 → configuração versionada → CSVs por execução → checksums/\allowbreak{}hashes → resumo com média e DP → interpretação limitada}



\section{Reprodução}



\begin{Verbatim}[fontsize=\small,breaklines=true,label=powershell]
$env:PYTHONPATH = "fases\01_validacao_conceitual"
.\.venv\Scripts\python.exe -m validacao.benchmark_operacoes --config <config.json> --output-dir <diretorio-novo>
.\.venv\Scripts\python.exe -m validacao.analise_benchmark_operacoes --evidence-dir <diretorio-novo>
\end{Verbatim}
